\chapter{Introduction}
\label{chap:intro}

\section{Motivation}

Railway transportation remains relevant even after its long history due to its efficiency, safety, and~sustainability. However, the~most limiting factor has always been the~use of~the required infrastructure. As such, the~critical problem of~utilizing the~railway effectively has always been scheduling. The complexity of~this problem is continuously increasing with the~size and~intricacy of~the railway system as well as the~demand to use it. 

Traditionally, scheduling was done manually by humans; only with the~advancements of~optimizing algorithms and~computational power have humans been aided by computers. This still produces a static schedule that might be unusable if a disruption occurs. This is where train dispatching comes into play where the~decisions are made in~real-time, managing delays and~conflicts as they happen. Real-time dispatching is often still only done by humans and~provides a challenge for~optimization algorithms as they need to find an acceptable solution in~a matter of~minutes or~even seconds.

\section{Goals}
The main goal of~this thesis is to develop a heuristic that is capable of~train dispatching in~real-time with interpretable results; a human should be able to understand the~structure of~the result to verify its validity.

The heuristic will be evaluated on~the \emph{DISPLIB}\@: Train dispatching benchmark library \cite{kloster2025displiblibrarytraindispatching} consisting of~train dispatching problem (TDP) instances from real-world data. This library of~instances was introduced for~an optimization competition.

A significant part of~this thesis is focused on~the aggregation of~the problem as presented in~the \emph{DISPLIB}. This is the~main innovation which allows us to use a top-down approach, reducing the~problem into selecting the~routes top-down and~then scheduling it with a more concise model, as opposed to the~bottom-up approach, building up the~solutions by sequentially routing and~resolving conflicts, which is used by the~winning solutions in~the \emph{DISPLIB} competition.

\section{Thesis structure}
\begin{minipage}{\textwidth}
This thesis is structured in~the following chapters with their contents:

\begin{enumerate}
	\item \hyperref[chap:intro]{\emph{Introduction:}} Outlines the~motivation, goals and~structure of~this thesis.
	%\hyperref[chap:lit_rev]
	\item \hyperref[chap:lit_rev]{\emph{Literature review:}} Review of~the pertinent literature separated into three parts, literature connected to the~\emph{DISPLIB}, the~most relevant resources inspiring the~heuristic or~providing key theoretical ideas, and~an overview of~selected state-of-the-art in~scheduling optimization.
	\item \hyperref[chap:theor]{\emph{Theoretical background:}} Basic theory necessary for~the description of~the problem. The~main focus is on~graph topology and~the use of~alternative and~disjunctive graphs for~scheduling models.
	\item \hyperref[chap:method]{\emph{Methodology:}} Detailed description of~the problem, aggregation, routing and~conflict resolution.
	\item \hyperref[chap:results]{\emph{Results:}} Presentation of~performance of~the algorithms described in~methodology on~the~\emph{DISPLIB} instances.
	\item \hyperref[chap:discuss]{\emph{Discussion:}} Discussion of the ideas and procedures presented in the methodology chapter and discussion of~results.
	\item \hyperref[chap:concl]{\emph{Conclusion:}} Highlight of~main ideas presented in~this thesis and overall evaluation, future work considerations.
\end{enumerate}
\end{minipage}


